% META: {"id": "TCOMPL-ATT-CO-005", "chapitre": "TCOMPL-TEMPS-ATTENTE", "type_objet": "corrige", "exercice_ref": "TCOMPL-ATT-EX-005", "capacites_codes": ["C1"], "sources_inspiration": [], "status": "generated"}
% BEGIN-VERIFY
% from sympy import Rational, symbols, summation, oo, simplify
% p = Rational(1, 6)
% def P(k):
%     return (1 - p) ** (k - 1) * p
% assert P(1) == Rational(1, 6)
% assert P(2) == Rational(5, 36)
% assert P(3) == Rational(25, 216)
% k = symbols('k', integer=True, positive=True)
% assert summation((1 - p) ** (k - 1) * p, (k, 1, oo)) == 1
% esperance = summation(k * (1 - p) ** (k - 1) * p, (k, 1, oo))
% assert simplify(esperance - 6) == 0
% END-VERIFY
\begin{corrige}{TCOMPL-ATT-EX-005}
\textbf{1.} On répète la même épreuve de Bernoulli de paramètre
$p=\tfrac16$, de façon indépendante, jusqu'au premier succès : $X$ suit la loi
géométrique de paramètre $\tfrac16$.

\textbf{2.} $P(X=k)=(1-p)^{k-1}p$, d'où $P(X=1)=\tfrac16$,
$P(X=2)=\tfrac56\times\tfrac16=\tfrac{5}{36}$ et
$P(X=3)=\left(\tfrac56\right)^2\tfrac16=\tfrac{25}{216}$.

\textbf{3.} $E(X)=\dfrac1p=6$ : il faut en moyenne six lancers pour obtenir un
premier six.
\end{corrige}
